\documentclass[a4paper, 12pt]{scrartcl}
\usepackage{comment}
\usepackage{graphicx}
\usepackage[utf8]{inputenc}
\usepackage[ngerman]{babel}
\usepackage[T1]{fontenc}
\usepackage{fancyhdr}
\usepackage{csquotes}
\pagestyle{fancy}
\usepackage[hidelinks]{hyperref}
\usepackage[center]{caption}
\usepackage[onehalfspacing]{setspace}
\usepackage[backend=bibtex, style=authoryear-ibid]{biblatex}
\usepackage{pgfplots}
\usepackage{listings}
%\usepackage[%
%    style=authoryear,
%    giveninits=true,
%    natbib=true,
%    maxbibnames=99,
%    uniquename=init
%]{biblatex}

\DeclareNameAlias{sortname}{family-given}
\renewcommand*{\multinamedelim}{\addcomma\space}
\renewcommand*{\finalnamedelim}{\addcomma\space}
\renewcommand*{\nameyeardelim}{\addcomma\space}
\setlength{\bibitemsep}{\baselineskip}
\addbibresource{lit.bib}

\begin{document}
    \renewcommand{\figurename}{Abb.}
	\fancyhf{}
	\thispagestyle{empty}
\begin{titlepage}
 \begin{center}
    \includegraphics[width=7.7cm]{logo_h-da_neu} \\ 
  \end{center}

  \begin{center}
    \vspace{0.1cm}
    \huge \textbf{Hochschule Darmstadt}\\
    \vspace{0.4cm}
    \LARGE -- Fachbereich Informatik --
  \end{center}

  \vfill
  \vfill

  \begin{center}
    \LARGE \textbf{Robustheit und Generalisierbarkeit in algorithmischen und Reinforcement Learning gestützten Lösungsansätzen: Eine Fallstudie mit Vier Gewinnt}
  \end{center} 
  
  \vfill
  \vfill

  \begin{center}
	\Large Abschlussarbeit zur Erlangung des akademischen Grades\\
	\vspace{0.3cm}
	\Large Bachelor of Science (B.Sc.)\\
  \end{center}

  \vfill

  \begin{center}
    \Large vorgelegt von\\
    \vspace{0.3cm}
    \Large \textbf{Leo Herrmann}\\
    \vspace{0.3cm}
    \normalsize Matrikelnummer: 1111455
  \end{center}
    \vfill
  \vfill

  \begin{center}
    \begin{tabular}{lll}
      Referentin:    & Prof. Dr. Elke Hergenröther \\
      Korreferent: & Adriatik Gashi				 \\
    \end{tabular}
  \end{center} 
\end{titlepage}

	\newpage
	
	\section{Kurzfassung}
	\newpage
	
	\tableofcontents
	\newpage

    \listoffigures
    \newpage

    \section*{Eigenständigkeitserklärung}
	Ich versichere hiermit, dass ich die vorliegende Arbeit selbständig verfasst und keine anderen als die im Literaturverzeichnis angegebenen Quellen benutzt habe.
Alle Stellen, die wörtlich oder sinngemäß aus veröffentlichten oder noch nicht veröffentlichten Quellen entnommen sind, sind als solche kenntlich gemacht.
Die Zeichnungen oder Abbildungen in dieser Arbeit sind von mir selbst erstellt worden oder mit einem entsprechenden Quellennachweis versehen.


\vspace{4cm}

Darmstadt, 07.03.2025

\hspace*{\fill}\begin{tabular}{@{}l@{}}\hline
\makebox[11cm]{}Leo Herrmann
\end{tabular}
	\newpage
 
	\setcounter{page}{1}
	\fancyfoot[C]{\thepage}
	
	\section{Einleitung}
	Fortschreitende Automatisierung durchdringt zahlreiche Bereiche der Gesellschaft, so zum Beispiel die Fertigungsindustrie, das Gesundheitswesen oder den Straßenverkehr. Zwei fundamentale Ansätze sind dabei symbolische Algorithmen und Machine Learning (ML). Die Einsatzbedingungen von Automatisierungssystemen unterscheiden sich häufig von den Bedingungen, unter denen sie entwickelt und getestet werden. Häufig müssen Systeme mit fehlerhaften oder veralteten Informationen arbeiten oder es treten Situationen ein, die bei der Konzipierung der Systeme nicht berücksichtigt werden können. Dabei sinkt die Leistungsfähigkeit dieser Systeme.

Bei Robustheit handelt es sich um eine Eigenschaft von Software, die beschreibt, wie gut sie unter genau solchen veränderten Bedingungen funktioniert. Robustheit ist daher besonders relevant für den Erfolg von Algorithmen und Modellen in der Praxis und wird in der Literatur ausgiebig untersucht \cite{Micskei.2012}\cite{Moos.2022}\cite{Ni.2021}.

Bei der Entscheidung zwischen symbolischen Algorithmen und ML-Verfahren sind entscheidende Faktoren in der Regel die Leistungsfähigkeit, Zuverlässigkeit und Komplexität der Lösung und die Nachvollziehbarkeit der Lösungsfindung (\cite{Dabas.2022}; \cite{Ferguson.January2019}, Kapitel 1.1.2; \cite{Humm.2020}, S. 12 f.; \cite{Früh.2022}, S. 5f.). Hierbei stellt sich die Frage, inwiefern bei der Entscheidung zwischen symbolischen und ML-basierten Lösungsansätzen auch Robustheit ein relevantes Kriterium ist, und falls ja, welche Art von Verfahren robuster ist.

Spiele eignen sich zur Untersuchung von Algorithmen und Modellen, weil sie reale Probleme auf kontrollierbare Umgebungen abstrahieren und gleichzeitig reproduzierbare und vergleichbare Messungen ermöglichen. Die Untersuchungen dieser Arbeit erfolgen am Beispiel des Brettspiels Vier Gewinnt, da aus früheren Untersuchungen ersichtlich wird, dass sich dafür sowohl symbolische Lösungsansätze als auch Verfahren aus dem ML-Teilbereich Reinforcement Learning eignen \cite{Qiu.2022}\cite{Sheoran.2022}\cite{Alderton.2019}\cite{Taylor.2024}\cite{Dabas.2022}\cite{Wäldchen.2022}.

Im Rahmen dieser Arbeit wird zunächst Grundlagenforschung zu verbreiteten symbolischen und Reinforcement Learning basierten Ansätzen betrieben. Dabei wird aus beiden Bereichen jeweils ein Verfahren ausgewählt, das sich zur Lösung von Vier Gewinnt und zur Untersuchung der Fragestellung eignet. Diese beiden Verfahren werden anschließend am Beispiel von Vier Gewinnt empirisch auf deren Robustheit untersucht. Dadurch sollen neue Erkenntnisse darüber gewonnen werden, wie Robustheit als Kriterium für Entscheidungen zwischen symbolischen Algorithmen und ML eingesetzt werden kann.

	
	\section{Grundlagen}
    In diesem Kapitel wird durch Literaturrecherche eine fundierte theoretische Basis geschaffen, auf die im weiteren Verlauf dieser Arbeit Bezug genommen wird. Zunächst wird Vier Gewinnt als zu lösendes Problem untersucht und eingeordnet. Anschließend folgt eine Auswahl von jeweils einem algorithmischen und einem Reinforcement Learning basiertem Lösungsansatz. Die Funktionsweise beider Lösungsmethoden wird erklärt. Außerdem werden bestehende Theorien und Definitionen zum Thema Robustheit und Generalisierbarkeit zusammengetragen. Sie bilden die Grundlage für die Szenarien und Bewertungskriterien in den Experimenten des Hauptteils.
	\subsection{Automatisierung von Spielen}
	% Hier geht es um Algorithmen, die Bäume durchsuchen, also erstmal.

\subsubsection{Minimax}

Minimax (auch Minmax) ist ein Algorithmus, der ausgehend von einem Knoten im Spielbaum die darauf folgenden Knoten bewertet und den Knoten mit der besten Bewertung zurückgibt. Bei der Bewertung wird davon ausgegangen, dass der Gegner ebenfalls den Zug wählt, der für sich am günstigsten ist. Das führt dazu, dass wenn die Bewertung anhand der Gewinnchancen erfolgt, auch tatsächlich die Gewinnchancen maximiert werden.

Um die Gewinnchancen zu ermitteln, müssen jedoch alle Knoten des Spielbaums untersucht werden. Die Laufzeit des Algorithmus steigt linear zur Anzahl der zu untersuchenden Knoten und damit bei konstanter Anzahl von Möglichkeiten pro Zug exponentiell zur Suchtiefe. Den gesamten Spielbaum zu durchsuchen, ist daher nur für wenig komplexe Spiele praktikabel. Für komplexere Spiele muss die Suchtiefe begrenzt und auf Heuristiken zurückgegriffen werden, damit die Bewertung in akzeptabler Zeit erfolgen kann\cite{Ferguson.January2019}\cite{Heineman.October2008}.

\subsubsection{AlphaBeta}

Bei AlphaBeta handelt es sich um eine Erweiterung von Minimax.

	\subsection{Vier Gewinnt}
    Vier Gewinnt ist ein Brettspiel, das aus einem 7 x 6 Spielfeld besteht. Die beiden Spieler werfen abwechselnd einen Spielstein in eine Spalte hinein, der in dieser Spalte bis zur untersten freien Position fällt. Es gewinnt der Spieler, der als erstes vier Spielsteine in einer horizontalen, vertikalen oder diagonalen Linie nebeneinander stehen hat \cite{MiltonBradleyCompany.1990}.

Bei Vier Gewinnt handelt es sich um ein kombinatorisches Nullsummenspiel für zwei Spieler. Kombinatorische Spiele weisen \glqq perfekte Information\grqq{} auf. Das bedeutet, dass alle Spieler zu jeder Zeit den gesamten Zustand des Spiels kennen. So ist es bei vielen Brettspielen der Fall. Kartenspiele hingegen besitzen diese Eigenschaft meistens nicht, weil jedem Spieler die Handkarten ihrer Gegenspieler unbekannt sind. Bei kombinatorischen Spielen sind außerdem keine Zufallselemente enthalten. Die einzige Herausforderung beim Spielen kombinatorischer Spiele besteht darin, unter einer Vielzahl von Entscheidungsoptionen diejenige auszuwählen, die den besten weiteren Spielverlauf verspricht (\cite{Bewersdorff.2018}, S. 96-100; \cite{Ferguson.January2019}, Kapitel 4.1).

Bei Zwei-Spieler-Nullsummenspielen verursacht der Gewinn eines Spielers zwangsläufig einen Verlust des anderen Spielers. Die beiden Spieler haben also entgegengesetzte Interessen (\cite{Bewersdorff.2018}, S. 100; \cite{Allis.1994}, S. 6). Das bedeutet, dass sich der Erfolg von verschiedenen Lösungsansätzen durch die durchschnittliche Gewinnrate im Spiel gegeneinander bewerten lässt. Bei Nullsummenspielen mit mehr als zwei Personen, kann es passieren, dass, wenn ein Spieler (bewusst oder versehentlich) nicht optimal spielt, ein zweiter Spieler davon profitiert, während ein dritter Spieler dadurch benachteiligt wird. Solche Wechselwirkungen sind bei Zwei-Personen-Nullsummenspielen ausgeschlossen (\cite{Bewersdorff.2018}, S. 113 ff.). Das macht die Messergebnisse im Hauptteil besser vergleichbar.

Nach Victor Allis lässt sich die Komplexität eines Spiels von strategiebasierten Zwei-Spieler-Nullsummenspielen durch ihre Zustandsraum- und Spielbaumkomplexität beschreiben. Die Zustandsraumkomplexität entspricht der Anzahl der verschiedenen möglichen Spielfeldkonfigurationen ab dem Start. Für ein Spiel kann dieser Wert oder zumindest dessen obere Schranke bestimmt werden, indem zunächst alle Konfigurationen des Spielfelds gezählt, dann Einschränkungen wie Regeln und Symmetrie berücksichtigt werden, und die Anzahl der illegalen und redundanten Zustände von der Anzahl aller möglichen Konfigurationen abgezogen wird (\cite{Allis.1994}, S. 158 f.).

Die Spielbaumkomplexität beschreibt die Anzahl der Blattknoten des Lösungsbaums. Der Spielbaum ist ein Baum, der die Zustände eines Spiels als Knoten und die Züge als Kanten darstellt (\cite{Bewersdorff.2018}, S. 102). Der Lösungsbaum beschreibt die Teilmenge des Spielbaums, der benötigt wird, um die Gewinnaussichten bei optimaler Spielweise beider Spieler zu berechnen. Die Spielbaumkomplexität lässt sich durch die durchschnittliche Spiellänge und der Anzahl der Entscheidungsmöglichkeiten pro Zug (entweder konstant oder abhängig vom Spielfortschritt) approximieren. Da in den meisten Spielen ein Zustand über mehrere Wege erreicht werden kann, fällt die Spielbaumkomplexität meist wesentlich größer aus als die Zustandsraumkomplexität (\cite{Allis.1994}, S. 159 ff.).

Die Spielbaumkomplexität ist maßgeblich für die praktische Berechenbarkeit einer starken Lösung. Für Tic Tac Toe wurde durch Allis eine obere Grenze für die Spielbaumkomplexität von 362880 ermittelt und eine starke Lösung lässt sich innerhalb von Sekundenbruchteilen berechnen \cite{Paul.2009}. Für Schach wird die Spielbaumkomplexität auf $10^{31}$ geschätzt und die Aussichten auf eine starke Lösung liegen noch in weiter Ferne \cite{Schaeffer.2007}.

Für Vier Gewinnt wurde eine durchschnittliche Spiellänge von 36 Zügen und eine durchschnittliche Anzahl von Entscheidungsmöglichkeiten (freie Spalten) von 4 ermittelt. Damit wurde die Spielbaumkomplexität auf $4^{36} \approx 10^{21}$ geschätzt (\cite{Allis.1994}, S. 163).

Verschiedene Lösungsverfahren von Vier Gewinnt sind bereits ausgiebig untersucht. Das Spiel wurde 1988 von James Dow Allen und Victor Allis unabhängig voneinander mit wissensbasierten Methoden schwach gelöst, was bedeutet, dass für die Anfangsposition eine optimale Strategie ermittelt wurde. Im Fall von Vier Gewinnt kann der Spieler, der den ersten Zug macht, bei optimaler Spielweise immer gewinnen \cite{Allen.2010}\cite{Allis.1988}.

1993 wurde das Spiel von John Tromp auch durch einen Brute-Force Ansatz stark gelöst. Bei dieser Lösung kam Alpha-Beta-Pruning zum Einsatz, um bei einer Zustandsraumkomplexität von 4.531.985.219.092 die optimalen Zugfolgen für beide Spieler zu berechnen. Das hat damals etwa 40.000 CPU-Stunden gedauert \cite{Tromp}.

Lösungen, die alle Möglichkeiten durchrechnen, um die optimale Entscheidung zu treffen, sind für den Einsatz in der Praxis aufgrund des hohen Rechenaufwands bei komplexeren Anwendungen auch heute noch selten praktikabel. Aus diesem Grund wird bevorzugt auf gute Heuristiken zurückgegriffen, die den Rechenaufwand minimieren, aber dennoch gute Ergebnisse liefern (\cite{Heineman.October2008}, Kapitel 7.6).

Untersuchungen haben gezeigt, dass sich sowohl regelbasierte Tree-Search-Algorithmen als auch verschiedene RL-Ansätze eignen, um sogenannte Agents zu entwickeln, die das Spiel selbstständig spielen \cite{Alderton.2019}\cite{Thill.2012}\cite{Wäldchen.2022}\cite{Taylor.2024}\cite{Sheoran.2022}\cite{Qiu.2022}. Wissensbasierte Methoden werden in dieser Arbeit nicht näher betrachtet, da ihre Ergebnisse stark an die jeweiligen Spielregeln gebunden und schwer zu verallgemeinern sind.


    \section{Konzept}
    In diesem Kapitel wird erklärt, wie eine Messumgebung aufgesetzt wurde, um die Robustheit der Lösungsverfahren empirisch zu bewerten. Diese Messumgebung ermöglicht es, zwei Agenten, die die zu untersuchenden Ansätze implementieren, das Spiel wiederholt gegeneinander spielen zu lassen. Die Spiele werden unter verschiedenen Szenarien durchgeführt, die jeweils auf ein in Kapitel \ref{robustheit} definierten Aspekt der Robustheit abzielen. Es werden die Gewinnraten gemessen, worüber Aussagen darüber getroffen werden können, welches der beiden Verfahren in den verschiedenen Szenarien stärker und damit robuster ist. Folgende Szenarien werden untersucht:

\begin{itemize}
	\item Unsicherheit bezüglich Aktionen: Die Aktionen der Agenten bestehen darin, dass sie den nächsten Spielstein in eine freie Spalte des Spielfelds hineinwerfen, die sie auf Grundlage des beobachteten Spielfeldzustands auswählen. Dieses Szenario führt Unsicherheit bezüglich Aktionen ein, indem für einen Agenten der Spielstein mit einer bestimmten Wahrscheinlichkeit nicht in die ausgewählte Spalte, sondern in eine zufällige freie Spalte fällt.
	
	\item Unsicherheit bezüglich Beobachtungen: Die Agenten wählen zwischen möglichen Aktionen basierend auf deren Beobachtungen des Spielfelds. In diesem Szenario erhalten die Agenten fehlerhafte Informationen über das Spielfeld. Jedes Feld besitzt dabei eine bestimmte Wahrscheinlichkeit mit der nicht dessen tatsächlicher Zustand (leer, besetzt durch Spieler 1, besetzt durch Spieler 2) erkannt wird, sondern ein zufälliger Zustand.
	
	Es ist anzumerken, dass hierbei aus dem MDP ein Partially Observable MDP (POMDP) wird, also ein MDP dessen Umgebung nur teilweise oder fehlerhaft beobachtbar ist. Der betroffene Agent kann nicht mit Sicherheit gesagt werden, in welchem Zustand er sich gerade befindet. Zusätzlich zum MDP enthält das POMDP ein Observation Modell O(s, o), das die Wahrscheinlichkeit beschreibt, eine Beobachtung o im Zustand s zu machen. POMDPs sind wesentlich komplizierter gezielt zu lösen, in der realen Welt jedoch wesentlich häufiger anzutreffen. Es gibt Optimierungen von MCTS und bestimmte RL-Methoden, über die POMDPs gezielt gelöst werden können, zum Beispiel indem für eine Entscheidung nicht nur der aktuelle Zustand, sondern die Historie der Zustände betrachtet wird (\cite{Russell.2020}, S. 588 ff.). Da es in dieser Arbeit darum geht, zu untersuchen, inwiefern grundsätzliche Eigenschaften von symbolischen Algorithmen und Reinforcement Learning Robustheit beeinflussen, werden diese gezielten Lösungen zur Vereinfachung nicht betrachtet. Beide Agenten gehen davon aus, dass es sich bei dem fehlerhaften Bild um den tatsächlichen Zustand des Spiels handelt, auch wenn der beobachtete Zustand gemäß der Spielregeln nicht erreicht werden könnte.
	
	\item Unsicherheit bezüglich Dynamik der Umgebung: Die Agenten erwarten ein bestimmtes Verhalten von der Umgebung, das durch die Spielregeln abgebildet wird. Der MCTS-Agent führt Simulationen anhand dieser Erwartungen aus und der RL-Agent wird auf Grundlage dieser Erwartungen trainiert. In diesem Szenario werden die Erwartungen an das Verhalten der Umgebung gebrochen, indem ein bestimmter Spieler mit einer bestimmten Wahrscheinlichkeit zwei Züge hintereinander durchführt.
	
	% \item Unsicherheit bezüglich Dynamik der Umgebung: Mit einer bestimmten Wahrscheinlichkeit führt der Gegenspieler nicht den Zug aus, den er für am besten hält, sondern einen zufälligen Zug. (Wurde rausgenommen, weil das quasi das gleiche ist wie Unsicherheit bezüglich Aktionen) nur von der anderen Perspektive.
\end{itemize}

In den verschiedenen Szenarien gelten die veränderten Bedingungen nur für jeweils einen Agenten. Wenn beide Agenten gleichzeitig von den veränderten Bedingungen betroffen wären, könnten nur Aussagen über die relative Robustheit zueinander getroffen werden. Dadurch dass die veränderten Bedingungen nur für jeweils ein Verfahren auf einmal gelten, kann genau gesagt werden, welches Verfahren wie stark betroffen ist.

Die Messungen werden nicht nur unter Szenarien mit veränderten Bedingungen durchgeführt, sondern auch in einer neutralen Umgebung ohne veränderte Bedingungen, um eine Grundlage für die Analyse der Ergebnisse zu bilden.

Um einen fairen Vergleich zu gewährleisten, müssen die Agenten in der neutralen Umgebung gleich stark sein. Dazu werden die Parameter der Agenten (Anzahl der Simulationen beim MCTS Agenten, Anzahl der durchlaufenen Self-Play-Trainingsepisonden beim RL-Agenten) so eingestellt, dass sie in der neutralen Umgebung im Spiel gegeneinander jeweils eine Gewinnrate von etwa 50\% aufweisen.

Außerdem ist es wichtig, dass die Agenten auf einem gewissen starken Level spielen. Denn bei Agenten, die ohnehin nicht stark spielen, wird es schwierig sein, in den verschiedenen Szenarien zur Untersuchung der Robustheit eine aussagekräftige Änderung in den Gewinnraten zu messen. Als Indikator für die Spielstärke der Agenten wird die durchschnittliche Spieldauer verwendet. Es wird angenommen, dass starke Agenten weniger Fehler machen und ihre Züge strategischer wählen, sodass komplexere Spielsituationen entstehen und sich die Entscheidung des Spiels hinauszögert.

Da bei Vier Gewinnt, wie in Kapitel \ref{vier-gewinnt} erwähnt, der Spieler, der den ersten Stein platziert, einen Vorteil hat, wechselt zu Beginn jedes Spiels das Recht, den ersten Zug zu machen.

    
    \section{Realisierung}
    \newpage
    
    \section{Ergebnisdiskussion}
    \newpage
    
    \section{Zusammenfassung und Ausblick}
    \newpage
 
	\section{Literaturverzeichnis}
	\nocite{*}
    \renewcommand{\refname}{}
	\printbibliography
 
	\newpage
	
\end{document}
